%% ============================================================
%%  diagrama_xx12.tex  —  Diagrama de flujo de xx12 (Erlang B)
%%  Uso: %% ============================================================
%%  diagrama_xx12.tex  —  Diagrama de flujo de xx12 (Erlang B)
%%  Uso: %% ============================================================
%%  diagrama_xx12.tex  —  Diagrama de flujo de xx12 (Erlang B)
%%  Uso: %% ============================================================
%%  diagrama_xx12.tex  —  Diagrama de flujo de xx12 (Erlang B)
%%  Uso: \input{diagrama_xx12}
%%
%%  Requiere en el preámbulo:
%%    \usepackage{tikz}
%%    \usepackage{xcolor}
%%    \usepackage{amsmath}
%%    \usetikzlibrary{shapes.geometric, shapes.misc,
%%      arrows.meta, shadows.blur, backgrounds,
%%      positioning, calc}
%% ============================================================

%% ── Paleta ──────────────────────────────────────────────────
\definecolor{fcbg}    {HTML}{F3F1EC}
\definecolor{fcBlue}  {HTML}{1C4E80}
\definecolor{fcTeal}  {HTML}{0A6B6B}
\definecolor{fcGreen} {HTML}{1E5C3A}
\definecolor{fcRed}   {HTML}{8B1A1A}
\definecolor{fcViolet}{HTML}{4A2080}
\definecolor{fcLine}  {HTML}{2C2C2C}

\definecolor{fcFBlue}  {HTML}{D6E4F0}
\definecolor{fcFTeal}  {HTML}{C8EAEA}
\definecolor{fcFGreen} {HTML}{C8E6C9}
\definecolor{fcFRed}   {HTML}{FADBD8}
\definecolor{fcFViolet}{HTML}{E8D5F5}

%% ── Estilos ─────────────────────────────────────────────────
\tikzset{
  fc base/.style={
    draw, align=center, inner sep=8pt,
    font=\small, line width=0.65pt,
    blur shadow={shadow blur steps=6,
                 shadow xshift=0.8pt, shadow yshift=-0.8pt,
                 shadow opacity=12},
  },
  fc term/.style={
    fc base, rounded rectangle,
    rounded rectangle arc length=180,
    minimum height=0.82cm, text width=3.8cm,
    font=\small\bfseries,
    draw=fcBlue, fill=fcFBlue, text=fcBlue,
  },
  fc proc/.style={
    fc base, rounded corners=4pt,
    minimum height=0.82cm, text width=5.2cm,
    draw=fcBlue, fill=fcFBlue, text=fcBlue,
  },
  fc dec/.style={
    fc base, diamond, aspect=2.8,
    minimum height=0.82cm, text width=4.4cm,
    inner sep=2pt, font=\small\bfseries,
    draw=fcTeal, fill=fcFTeal, text=fcTeal,
  },
  fc loop/.style={
    fc base, rounded corners=4pt,
    minimum height=0.82cm, text width=5.2cm,
    font=\small\bfseries,
    draw=fcGreen, fill=fcFGreen, text=fcGreen,
  },
  fc recur/.style={
    fc base, rounded corners=4pt,
    minimum height=1.1cm, text width=5.2cm,
    draw=fcViolet, fill=fcFViolet, text=fcViolet,
  },
  fc err/.style={
    fc base, rounded corners=4pt,
    minimum height=0.82cm, text width=4.0cm,
    font=\small\bfseries,
    draw=fcRed, fill=fcFRed, text=fcRed,
  },
  fc arr/.style={
    ->, >=Stealth, line width=0.9pt, color=fcLine,
    shorten >=2pt, shorten <=1pt,
  },
  fc back/.style={
    ->, >=Stealth, line width=0.85pt,
    color=fcGreen, dashed,
    shorten >=2pt, shorten <=1pt,
  },
  fc lbl/.style={
    font=\footnotesize\bfseries,
    fill=fcbg, inner sep=1.5pt, rounded corners=1.5pt,
  },
}

%% ── Figura ──────────────────────────────────────────────────
\begin{tikzpicture}[node distance=0.65cm and 2.6cm]

%% ── Columna central ─────────────────────────────────────────
\node[fc term] (start)
  {Inicio \;\texttt{xx12(A,\,k,\,mostrar)}};

\node[fc dec, below=of start] (nargin)
  {\texttt{nargin} $= 2$?};

\node[fc proc, below=of nargin] (setmos)
  {\texttt{mostrar} $\leftarrow 0$};

\node[fc dec, below=0.8cm of setmos] (valid)
  {Parámetros válidos?};

\node[fc proc, below=0.8cm of valid] (init)
  {$p_k \leftarrow 1$,\quad $n \leftarrow 1$};

\node[fc loop, below=of init] (loopcond)
  {$n \leq k$?};

\node[fc recur, below=of loopcond] (recurse)
  {$\displaystyle p_k \leftarrow
    \frac{A\,p_k}{n + A\,p_k}$\quad
   $n \leftarrow n+1$};

\node[fc dec, below=1.1cm of recurse] (showdec)
  {\texttt{mostrar} $= 1$?};

\node[fc proc, below=of showdec] (print)
  {Imprimir $A$,\; $k$,\; $p_k$};

\node[fc term, below=of print] (stop)
  {Retornar $p_k$\; — Fin};

%% ── Laterales ───────────────────────────────────────────────
\node[fc err, right=of valid, xshift=0.8cm] (err)
  {\texttt{error(...)}\\Fin};

%% ── Flechas centrales ───────────────────────────────────────
\draw[fc arr] (start)    -- (nargin);
\draw[fc arr] (nargin)   -- node[fc lbl,left]{Sí} (setmos);
\draw[fc arr] (setmos)   -- (valid);
\draw[fc arr] (valid)    -- node[fc lbl,left]{Sí} (init);
\draw[fc arr] (init)     -- (loopcond);
\draw[fc arr] (loopcond) -- node[fc lbl,left]{Sí} (recurse);
\draw[fc arr] (showdec)  -- node[fc lbl,left]{Sí} (print);
\draw[fc arr] (print)    -- (stop);

%% nargin No → salta setmos, va directo a valid
\draw[fc arr]
  (nargin.east) -- ++(1.6,0)
  |- node[fc lbl, pos=0.18]{No} (valid.east);

%% valid No → error
\draw[fc arr]
  (valid.east) -- node[fc lbl,above]{} (err.west);

%% loopcond No → showdec (derecha)
\draw[fc arr]
  (loopcond.east) -- ++(1.8,0)
  |- node[fc lbl, pos=0.18]{No} (showdec.east);

%% showdec No → stop (derecha)
\draw[fc arr]
  (showdec.west) -- ++(-1.8,0)
  |- node[fc lbl, pos=0.18]{No} (stop.west);

%% Retorno del bucle (izquierda)
\draw[fc back]
  (recurse.west) -- ++(-1.8,0)
  |- (loopcond.west);

%% ── Fondo ───────────────────────────────────────────────────


\node[font=\tiny\bfseries, color=fcGreen, rotate=90, anchor=south]
  at ($(init.north west)+(-1.65,-0.3)$)
  {for $n=1\ldots k$};

%% ── Título ──────────────────────────────────────────────────
\node[font=\normalsize\bfseries, color=fcBlue, anchor=south]
  at ($(start.north)+(0,0.52)$)
  {Diagrama de flujo — \texttt{xx12} \;($M/M/k/k$, Erlang B)};

\end{tikzpicture}

%%
%%  Requiere en el preámbulo:
%%    \usepackage{tikz}
%%    \usepackage{xcolor}
%%    \usepackage{amsmath}
%%    \usetikzlibrary{shapes.geometric, shapes.misc,
%%      arrows.meta, shadows.blur, backgrounds,
%%      positioning, calc}
%% ============================================================

%% ── Paleta ──────────────────────────────────────────────────
\definecolor{fcbg}    {HTML}{F3F1EC}
\definecolor{fcBlue}  {HTML}{1C4E80}
\definecolor{fcTeal}  {HTML}{0A6B6B}
\definecolor{fcGreen} {HTML}{1E5C3A}
\definecolor{fcRed}   {HTML}{8B1A1A}
\definecolor{fcViolet}{HTML}{4A2080}
\definecolor{fcLine}  {HTML}{2C2C2C}

\definecolor{fcFBlue}  {HTML}{D6E4F0}
\definecolor{fcFTeal}  {HTML}{C8EAEA}
\definecolor{fcFGreen} {HTML}{C8E6C9}
\definecolor{fcFRed}   {HTML}{FADBD8}
\definecolor{fcFViolet}{HTML}{E8D5F5}

%% ── Estilos ─────────────────────────────────────────────────
\tikzset{
  fc base/.style={
    draw, align=center, inner sep=8pt,
    font=\small, line width=0.65pt,
    blur shadow={shadow blur steps=6,
                 shadow xshift=0.8pt, shadow yshift=-0.8pt,
                 shadow opacity=12},
  },
  fc term/.style={
    fc base, rounded rectangle,
    rounded rectangle arc length=180,
    minimum height=0.82cm, text width=3.8cm,
    font=\small\bfseries,
    draw=fcBlue, fill=fcFBlue, text=fcBlue,
  },
  fc proc/.style={
    fc base, rounded corners=4pt,
    minimum height=0.82cm, text width=5.2cm,
    draw=fcBlue, fill=fcFBlue, text=fcBlue,
  },
  fc dec/.style={
    fc base, diamond, aspect=2.8,
    minimum height=0.82cm, text width=4.4cm,
    inner sep=2pt, font=\small\bfseries,
    draw=fcTeal, fill=fcFTeal, text=fcTeal,
  },
  fc loop/.style={
    fc base, rounded corners=4pt,
    minimum height=0.82cm, text width=5.2cm,
    font=\small\bfseries,
    draw=fcGreen, fill=fcFGreen, text=fcGreen,
  },
  fc recur/.style={
    fc base, rounded corners=4pt,
    minimum height=1.1cm, text width=5.2cm,
    draw=fcViolet, fill=fcFViolet, text=fcViolet,
  },
  fc err/.style={
    fc base, rounded corners=4pt,
    minimum height=0.82cm, text width=4.0cm,
    font=\small\bfseries,
    draw=fcRed, fill=fcFRed, text=fcRed,
  },
  fc arr/.style={
    ->, >=Stealth, line width=0.9pt, color=fcLine,
    shorten >=2pt, shorten <=1pt,
  },
  fc back/.style={
    ->, >=Stealth, line width=0.85pt,
    color=fcGreen, dashed,
    shorten >=2pt, shorten <=1pt,
  },
  fc lbl/.style={
    font=\footnotesize\bfseries,
    fill=fcbg, inner sep=1.5pt, rounded corners=1.5pt,
  },
}

%% ── Figura ──────────────────────────────────────────────────
\begin{tikzpicture}[node distance=0.65cm and 2.6cm]

%% ── Columna central ─────────────────────────────────────────
\node[fc term] (start)
  {Inicio \;\texttt{xx12(A,\,k,\,mostrar)}};

\node[fc dec, below=of start] (nargin)
  {\texttt{nargin} $= 2$?};

\node[fc proc, below=of nargin] (setmos)
  {\texttt{mostrar} $\leftarrow 0$};

\node[fc dec, below=0.8cm of setmos] (valid)
  {Parámetros válidos?};

\node[fc proc, below=0.8cm of valid] (init)
  {$p_k \leftarrow 1$,\quad $n \leftarrow 1$};

\node[fc loop, below=of init] (loopcond)
  {$n \leq k$?};

\node[fc recur, below=of loopcond] (recurse)
  {$\displaystyle p_k \leftarrow
    \frac{A\,p_k}{n + A\,p_k}$\quad
   $n \leftarrow n+1$};

\node[fc dec, below=1.1cm of recurse] (showdec)
  {\texttt{mostrar} $= 1$?};

\node[fc proc, below=of showdec] (print)
  {Imprimir $A$,\; $k$,\; $p_k$};

\node[fc term, below=of print] (stop)
  {Retornar $p_k$\; — Fin};

%% ── Laterales ───────────────────────────────────────────────
\node[fc err, right=of valid, xshift=0.8cm] (err)
  {\texttt{error(...)}\\Fin};

%% ── Flechas centrales ───────────────────────────────────────
\draw[fc arr] (start)    -- (nargin);
\draw[fc arr] (nargin)   -- node[fc lbl,left]{Sí} (setmos);
\draw[fc arr] (setmos)   -- (valid);
\draw[fc arr] (valid)    -- node[fc lbl,left]{Sí} (init);
\draw[fc arr] (init)     -- (loopcond);
\draw[fc arr] (loopcond) -- node[fc lbl,left]{Sí} (recurse);
\draw[fc arr] (showdec)  -- node[fc lbl,left]{Sí} (print);
\draw[fc arr] (print)    -- (stop);

%% nargin No → salta setmos, va directo a valid
\draw[fc arr]
  (nargin.east) -- ++(1.6,0)
  |- node[fc lbl, pos=0.18]{No} (valid.east);

%% valid No → error
\draw[fc arr]
  (valid.east) -- node[fc lbl,above]{} (err.west);

%% loopcond No → showdec (derecha)
\draw[fc arr]
  (loopcond.east) -- ++(1.8,0)
  |- node[fc lbl, pos=0.18]{No} (showdec.east);

%% showdec No → stop (derecha)
\draw[fc arr]
  (showdec.west) -- ++(-1.8,0)
  |- node[fc lbl, pos=0.18]{No} (stop.west);

%% Retorno del bucle (izquierda)
\draw[fc back]
  (recurse.west) -- ++(-1.8,0)
  |- (loopcond.west);

%% ── Fondo ───────────────────────────────────────────────────


\node[font=\tiny\bfseries, color=fcGreen, rotate=90, anchor=south]
  at ($(init.north west)+(-1.65,-0.3)$)
  {for $n=1\ldots k$};

%% ── Título ──────────────────────────────────────────────────
\node[font=\normalsize\bfseries, color=fcBlue, anchor=south]
  at ($(start.north)+(0,0.52)$)
  {Diagrama de flujo — \texttt{xx12} \;($M/M/k/k$, Erlang B)};

\end{tikzpicture}

%%
%%  Requiere en el preámbulo:
%%    \usepackage{tikz}
%%    \usepackage{xcolor}
%%    \usepackage{amsmath}
%%    \usetikzlibrary{shapes.geometric, shapes.misc,
%%      arrows.meta, shadows.blur, backgrounds,
%%      positioning, calc}
%% ============================================================

%% ── Paleta ──────────────────────────────────────────────────
\definecolor{fcbg}    {HTML}{F3F1EC}
\definecolor{fcBlue}  {HTML}{1C4E80}
\definecolor{fcTeal}  {HTML}{0A6B6B}
\definecolor{fcGreen} {HTML}{1E5C3A}
\definecolor{fcRed}   {HTML}{8B1A1A}
\definecolor{fcViolet}{HTML}{4A2080}
\definecolor{fcLine}  {HTML}{2C2C2C}

\definecolor{fcFBlue}  {HTML}{D6E4F0}
\definecolor{fcFTeal}  {HTML}{C8EAEA}
\definecolor{fcFGreen} {HTML}{C8E6C9}
\definecolor{fcFRed}   {HTML}{FADBD8}
\definecolor{fcFViolet}{HTML}{E8D5F5}

%% ── Estilos ─────────────────────────────────────────────────
\tikzset{
  fc base/.style={
    draw, align=center, inner sep=8pt,
    font=\small, line width=0.65pt,
    blur shadow={shadow blur steps=6,
                 shadow xshift=0.8pt, shadow yshift=-0.8pt,
                 shadow opacity=12},
  },
  fc term/.style={
    fc base, rounded rectangle,
    rounded rectangle arc length=180,
    minimum height=0.82cm, text width=3.8cm,
    font=\small\bfseries,
    draw=fcBlue, fill=fcFBlue, text=fcBlue,
  },
  fc proc/.style={
    fc base, rounded corners=4pt,
    minimum height=0.82cm, text width=5.2cm,
    draw=fcBlue, fill=fcFBlue, text=fcBlue,
  },
  fc dec/.style={
    fc base, diamond, aspect=2.8,
    minimum height=0.82cm, text width=4.4cm,
    inner sep=2pt, font=\small\bfseries,
    draw=fcTeal, fill=fcFTeal, text=fcTeal,
  },
  fc loop/.style={
    fc base, rounded corners=4pt,
    minimum height=0.82cm, text width=5.2cm,
    font=\small\bfseries,
    draw=fcGreen, fill=fcFGreen, text=fcGreen,
  },
  fc recur/.style={
    fc base, rounded corners=4pt,
    minimum height=1.1cm, text width=5.2cm,
    draw=fcViolet, fill=fcFViolet, text=fcViolet,
  },
  fc err/.style={
    fc base, rounded corners=4pt,
    minimum height=0.82cm, text width=4.0cm,
    font=\small\bfseries,
    draw=fcRed, fill=fcFRed, text=fcRed,
  },
  fc arr/.style={
    ->, >=Stealth, line width=0.9pt, color=fcLine,
    shorten >=2pt, shorten <=1pt,
  },
  fc back/.style={
    ->, >=Stealth, line width=0.85pt,
    color=fcGreen, dashed,
    shorten >=2pt, shorten <=1pt,
  },
  fc lbl/.style={
    font=\footnotesize\bfseries,
    fill=fcbg, inner sep=1.5pt, rounded corners=1.5pt,
  },
}

%% ── Figura ──────────────────────────────────────────────────
\begin{tikzpicture}[node distance=0.65cm and 2.6cm]

%% ── Columna central ─────────────────────────────────────────
\node[fc term] (start)
  {Inicio \;\texttt{xx12(A,\,k,\,mostrar)}};

\node[fc dec, below=of start] (nargin)
  {\texttt{nargin} $= 2$?};

\node[fc proc, below=of nargin] (setmos)
  {\texttt{mostrar} $\leftarrow 0$};

\node[fc dec, below=0.8cm of setmos] (valid)
  {Parámetros válidos?};

\node[fc proc, below=0.8cm of valid] (init)
  {$p_k \leftarrow 1$,\quad $n \leftarrow 1$};

\node[fc loop, below=of init] (loopcond)
  {$n \leq k$?};

\node[fc recur, below=of loopcond] (recurse)
  {$\displaystyle p_k \leftarrow
    \frac{A\,p_k}{n + A\,p_k}$\quad
   $n \leftarrow n+1$};

\node[fc dec, below=1.1cm of recurse] (showdec)
  {\texttt{mostrar} $= 1$?};

\node[fc proc, below=of showdec] (print)
  {Imprimir $A$,\; $k$,\; $p_k$};

\node[fc term, below=of print] (stop)
  {Retornar $p_k$\; — Fin};

%% ── Laterales ───────────────────────────────────────────────
\node[fc err, right=of valid, xshift=0.8cm] (err)
  {\texttt{error(...)}\\Fin};

%% ── Flechas centrales ───────────────────────────────────────
\draw[fc arr] (start)    -- (nargin);
\draw[fc arr] (nargin)   -- node[fc lbl,left]{Sí} (setmos);
\draw[fc arr] (setmos)   -- (valid);
\draw[fc arr] (valid)    -- node[fc lbl,left]{Sí} (init);
\draw[fc arr] (init)     -- (loopcond);
\draw[fc arr] (loopcond) -- node[fc lbl,left]{Sí} (recurse);
\draw[fc arr] (showdec)  -- node[fc lbl,left]{Sí} (print);
\draw[fc arr] (print)    -- (stop);

%% nargin No → salta setmos, va directo a valid
\draw[fc arr]
  (nargin.east) -- ++(1.6,0)
  |- node[fc lbl, pos=0.18]{No} (valid.east);

%% valid No → error
\draw[fc arr]
  (valid.east) -- node[fc lbl,above]{} (err.west);

%% loopcond No → showdec (derecha)
\draw[fc arr]
  (loopcond.east) -- ++(1.8,0)
  |- node[fc lbl, pos=0.18]{No} (showdec.east);

%% showdec No → stop (derecha)
\draw[fc arr]
  (showdec.west) -- ++(-1.8,0)
  |- node[fc lbl, pos=0.18]{No} (stop.west);

%% Retorno del bucle (izquierda)
\draw[fc back]
  (recurse.west) -- ++(-1.8,0)
  |- (loopcond.west);

%% ── Fondo ───────────────────────────────────────────────────


\node[font=\tiny\bfseries, color=fcGreen, rotate=90, anchor=south]
  at ($(init.north west)+(-1.65,-0.3)$)
  {for $n=1\ldots k$};

%% ── Título ──────────────────────────────────────────────────
\node[font=\normalsize\bfseries, color=fcBlue, anchor=south]
  at ($(start.north)+(0,0.52)$)
  {Diagrama de flujo — \texttt{xx12} \;($M/M/k/k$, Erlang B)};

\end{tikzpicture}

%%
%%  Requiere en el preámbulo:
%%    \usepackage{tikz}
%%    \usepackage{xcolor}
%%    \usepackage{amsmath}
%%    \usetikzlibrary{shapes.geometric, shapes.misc,
%%      arrows.meta, shadows.blur, backgrounds,
%%      positioning, calc}
%% ============================================================

%% ── Paleta ──────────────────────────────────────────────────
\definecolor{fcbg}    {HTML}{F3F1EC}
\definecolor{fcBlue}  {HTML}{1C4E80}
\definecolor{fcTeal}  {HTML}{0A6B6B}
\definecolor{fcGreen} {HTML}{1E5C3A}
\definecolor{fcRed}   {HTML}{8B1A1A}
\definecolor{fcViolet}{HTML}{4A2080}
\definecolor{fcLine}  {HTML}{2C2C2C}

\definecolor{fcFBlue}  {HTML}{D6E4F0}
\definecolor{fcFTeal}  {HTML}{C8EAEA}
\definecolor{fcFGreen} {HTML}{C8E6C9}
\definecolor{fcFRed}   {HTML}{FADBD8}
\definecolor{fcFViolet}{HTML}{E8D5F5}

%% ── Estilos ─────────────────────────────────────────────────
\tikzset{
  fc base/.style={
    draw, align=center, inner sep=8pt,
    font=\small, line width=0.65pt,
    blur shadow={shadow blur steps=6,
                 shadow xshift=0.8pt, shadow yshift=-0.8pt,
                 shadow opacity=12},
  },
  fc term/.style={
    fc base, rounded rectangle,
    rounded rectangle arc length=180,
    minimum height=0.82cm, text width=3.8cm,
    font=\small\bfseries,
    draw=fcBlue, fill=fcFBlue, text=fcBlue,
  },
  fc proc/.style={
    fc base, rounded corners=4pt,
    minimum height=0.82cm, text width=5.2cm,
    draw=fcBlue, fill=fcFBlue, text=fcBlue,
  },
  fc dec/.style={
    fc base, diamond, aspect=2.8,
    minimum height=0.82cm, text width=4.4cm,
    inner sep=2pt, font=\small\bfseries,
    draw=fcTeal, fill=fcFTeal, text=fcTeal,
  },
  fc loop/.style={
    fc base, rounded corners=4pt,
    minimum height=0.82cm, text width=5.2cm,
    font=\small\bfseries,
    draw=fcGreen, fill=fcFGreen, text=fcGreen,
  },
  fc recur/.style={
    fc base, rounded corners=4pt,
    minimum height=1.1cm, text width=5.2cm,
    draw=fcViolet, fill=fcFViolet, text=fcViolet,
  },
  fc err/.style={
    fc base, rounded corners=4pt,
    minimum height=0.82cm, text width=4.0cm,
    font=\small\bfseries,
    draw=fcRed, fill=fcFRed, text=fcRed,
  },
  fc arr/.style={
    ->, >=Stealth, line width=0.9pt, color=fcLine,
    shorten >=2pt, shorten <=1pt,
  },
  fc back/.style={
    ->, >=Stealth, line width=0.85pt,
    color=fcGreen, dashed,
    shorten >=2pt, shorten <=1pt,
  },
  fc lbl/.style={
    font=\footnotesize\bfseries,
    fill=fcbg, inner sep=1.5pt, rounded corners=1.5pt,
  },
}

%% ── Figura ──────────────────────────────────────────────────
\begin{tikzpicture}[node distance=0.65cm and 2.6cm]

%% ── Columna central ─────────────────────────────────────────
\node[fc term] (start)
  {Inicio \;\texttt{xx12(A,\,k,\,mostrar)}};

\node[fc dec, below=of start] (nargin)
  {\texttt{nargin} $= 2$?};

\node[fc proc, below=of nargin] (setmos)
  {\texttt{mostrar} $\leftarrow 0$};

\node[fc dec, below=0.8cm of setmos] (valid)
  {Parámetros válidos?};

\node[fc proc, below=0.8cm of valid] (init)
  {$p_k \leftarrow 1$,\quad $n \leftarrow 1$};

\node[fc loop, below=of init] (loopcond)
  {$n \leq k$?};

\node[fc recur, below=of loopcond] (recurse)
  {$\displaystyle p_k \leftarrow
    \frac{A\,p_k}{n + A\,p_k}$\quad
   $n \leftarrow n+1$};

\node[fc dec, below=1.1cm of recurse] (showdec)
  {\texttt{mostrar} $= 1$?};

\node[fc proc, below=of showdec] (print)
  {Imprimir $A$,\; $k$,\; $p_k$};

\node[fc term, below=of print] (stop)
  {Retornar $p_k$\; — Fin};

%% ── Laterales ───────────────────────────────────────────────
\node[fc err, right=of valid, xshift=0.8cm] (err)
  {\texttt{error(...)}\\Fin};

%% ── Flechas centrales ───────────────────────────────────────
\draw[fc arr] (start)    -- (nargin);
\draw[fc arr] (nargin)   -- node[fc lbl,left]{Sí} (setmos);
\draw[fc arr] (setmos)   -- (valid);
\draw[fc arr] (valid)    -- node[fc lbl,left]{Sí} (init);
\draw[fc arr] (init)     -- (loopcond);
\draw[fc arr] (loopcond) -- node[fc lbl,left]{Sí} (recurse);
\draw[fc arr] (showdec)  -- node[fc lbl,left]{Sí} (print);
\draw[fc arr] (print)    -- (stop);

%% nargin No → salta setmos, va directo a valid
\draw[fc arr]
  (nargin.east) -- ++(1.6,0)
  |- node[fc lbl, pos=0.18]{No} (valid.east);

%% valid No → error
\draw[fc arr]
  (valid.east) -- node[fc lbl,above]{} (err.west);

%% loopcond No → showdec (derecha)
\draw[fc arr]
  (loopcond.east) -- ++(1.8,0)
  |- node[fc lbl, pos=0.18]{No} (showdec.east);

%% showdec No → stop (derecha)
\draw[fc arr]
  (showdec.west) -- ++(-1.8,0)
  |- node[fc lbl, pos=0.18]{No} (stop.west);

%% Retorno del bucle (izquierda)
\draw[fc back]
  (recurse.west) -- ++(-1.8,0)
  |- (loopcond.west);

%% ── Fondo ───────────────────────────────────────────────────


\node[font=\tiny\bfseries, color=fcGreen, rotate=90, anchor=south]
  at ($(init.north west)+(-1.65,-0.3)$)
  {for $n=1\ldots k$};

%% ── Título ──────────────────────────────────────────────────
\node[font=\normalsize\bfseries, color=fcBlue, anchor=south]
  at ($(start.north)+(0,0.52)$)
  {Diagrama de flujo — \texttt{xx12} \;($M/M/k/k$, Erlang B)};

\end{tikzpicture}
